\chapter{Validation and limits}\label{ch:validation}

\section{Earlier control-mapping revision checks}
The earlier control-mapping browser checks cover explicit sample/control matching, negative matrix versus endogenous reference, retained aliquot
identity, valid and invalid extraction dates, a positive event after a zero-event baseline, incomplete baselines, mixed
instrument cohorts and the review-only laboratory scaffold. The generator rebuilds the synthetic runs, screenshots,
control-chart ZIP and code extracts from the revised HTML. Its identifier check examines textual metadata, including
Map values; coincidental simulated numeric measurements are not treated as identifiers.

The original release reported cell-by-cell agreement with vendor exports. That historical result is not a new validation
claim for this revision. The correction does not intentionally change the vendor export layout. Laboratory acceptance,
historical method applicability and the meaning of ambiguous aliquot labels still require verification.

Current Roche-reconstruction and offline performance checks are described in Chapter~\ref{ch:roche}. The earlier gallery was retained; the complete browser screenshot suite was not rerun for this update.

\section{Known limits}
\begin{itemize}
\item The page reads what the instrument stored. It cannot verify real well temperatures (a temperature-verification plate does)
      or light output (a fluorescence reference plate does).
\item Lamp hours, plate loading and the error log of the LightCycler are not in the \texttt{.ixo}; they are in an Instrument Problem
      Report. QuantStudio LED hours and drawer events are not in the \texttt{.eds}.
\item The QuantStudio \texttt{messages.log} is not a documented interface; a firmware update may rename lines. Version changes are
      marked on every chart.
\item Control limits need about 20 runs per instrument and program; before that the tiles stay grey.
\item Operator charts need an individual operator per run; a shared account groups all runs under one name.
\item The synthetic data set of this document is illustrative. Its trends and faults were chosen to show the charts, not to
      describe the behaviour of any instrument.
\end{itemize}
